\documentclass[11pt,twocolumn,a4paper]{article}
\usepackage{amsmath,amssymb,amsthm,graphicx,booktabs,hyperref,geometry,siunitx}
\usepackage{float}
\usepackage{natbib}
\usepackage{algorithm}
\usepackage{algorithmic}
\usepackage{listings}
\usepackage{xcolor}
\usepackage{caption}
\usepackage{microtype}
\usepackage{lmodern}
\usepackage[version=4]{mhchem}
\geometry{margin=2.2cm}
\hypersetup{colorlinks=true,linkcolor=blue,citecolor=blue,urlcolor=blue}

%% --- Shared monograph notation (Papers 1--2) -------------------------
\newcommand{\Sk}{S_{k}}
\newcommand{\St}{S_{t}}
\newcommand{\Se}{S_{e}}
\newcommand{\Sspace}{\mathcal{S}}
\newcommand{\Scoord}{s}
\newcommand{\Sexp}{\xi}
\newcommand{\addr}{\mathrm{addr}}
\newcommand{\trit}{\mathrm{trit}}
\newcommand{\eval}{\mathrm{eval}}
\newcommand{\Recv}{\mathcal{R}}
\newcommand{\Tree}{\mathcal{T}}
\newcommand{\Leaf}{\Lambda}
\newcommand{\dcat}{d_{\mathrm{cat}}}
\newcommand{\Depth}{\mathcal{M}}
\newcommand{\Part}{\mathcal{P}}
\newcommand{\dC}{d_{\mathcal{C}}}
\newcommand{\nufloor}{\nu_{\mathrm{floor}}}
\newcommand{\Res}{\mathcal{A}}
\newcommand{\Sch}{\Sigma}
\newcommand{\Cache}{\mathcal{K}}

\theoremstyle{plain}
\newtheorem{theorem}{Theorem}
\newtheorem{lemma}[theorem]{Lemma}
\newtheorem{corollary}[theorem]{Corollary}
\newtheorem{proposition}[theorem]{Proposition}
\theoremstyle{definition}
\newtheorem{definition}{Definition}
\newtheorem{construction}[definition]{Construction}
\theoremstyle{remark}
\newtheorem{remark}{Remark}
\newtheorem{example}[remark]{Example}

\title{\textbf{The Empty Dictionary for Cytochrome P450:\\
Constant Resident State, Query Without Entries,\\
and a Correction to the Storage Claim of Paper~2}}

\author{Levinthal Monograph Series --- Paper 16}
\date{}

\begin{document}
\maketitle

\begin{abstract}
Every preceding paper of this monograph computes \emph{within} a cytochrome
P450 evaluation: the sequence manifold, the fold, the catalytic cycle, the
isoform taxonomy, the spectra, the pharmacogenomics. None asks whether P450
knowledge must be \emph{stored} at all. This paper answers that question.
We instantiate the Empty Dictionary result of \citet{sachikonye2025c}---whose
worked case is SOD1---on cytochrome P450, in the monograph's own notation,
and we establish two things by measurement rather than assertion. First, the
\emph{resident state} of the P450 addressing scheme is $O(1)$ in the number
of addressable objects: it is exactly the twenty rows of the residue
coordinate table plus an encoding rule, and it is measured at \textbf{561
bytes}, invariant across corpora from a single sequence to $10^{9}$. Second,
queries---Identify, Similar, Predict---are answered for sequences that were
never stored, by evaluating a coordinate map, with the resident state
provably unmodified by the query batch. Two controls make these claims
non-vacuous: a conventional exact-match dictionary answers $0/6$ of the same
queries, and $200$ shuffled coordinate tables (same size, same address
algebra, wrong assignment) collapse the graded physicochemical correlation
from $\rho = 0.382$ to a mean of $-0.004$ (empirical $p < 0.005$ over $190$
residue pairs). We then correct Paper~2 of this monograph, which invokes the
Empty Dictionary Principle by name and, three paragraphs later, reports
storage of $O(N \cdot k)$ trits per sequence---${\sim}350$~MB for
$4 \times 10^{5}$ sequences. Both of Paper~2's arithmetic figures reproduce
($7.13$ vs.\ ${\sim}7$~kbits; $356.6$ vs.\ ${\sim}350$~MB), so the
disagreement is not arithmetical but categorial: that figure measures an
optional materialised cache, not the resident state. We show the cache is not
load-bearing---answers are byte-for-byte identical with it present and
deleted---and conclude that Paper~2 describes a \emph{compressed} dictionary
while calling it an empty one.
\end{abstract}

\tableofcontents
\newpage

%% ============================================================
\section{Introduction}
%% ============================================================

\subsection{The question the monograph has not asked}

The sixteen papers preceding this one share a presupposition so uniform that
it has never been stated. Each takes some cytochrome P450 object---a
sequence, a fold, a catalytic intermediate, an allele---and computes a
property of it. The presupposition is that the object is \emph{available}:
that it has been read from somewhere, that a representation of it exists
prior to the computation.

Paper~14 comes closest to examining this. It asks whether a corrupted P450
database can be recovered from partial ternary addresses, and answers
quantitatively: a $70\,\%$-complete $k=6$ address carries ${\approx}6.7$
bits against the $5.8$ needed for unique isoform identification. But the
question is posed \emph{inside} the presupposition. A database exists; is it
recoverable? The prior question---must there be a database?---is not asked.

This paper asks it, and the answer is no.

\subsection{What is being claimed, and what is not}

The claim is easy to overstate, so it is worth fixing its shape before any
machinery appears.

We do \emph{not} claim that P450 sequences carry no information, that
experimental data is dispensable, or that a laboratory can be replaced by a
coordinate table. The sequences are real, and someone must measure them.

What we claim is narrower and, we think, more interesting: the
\emph{persistent state} required to make P450 queries answerable does not
grow with the number of P450 objects that can be queried. A conventional
database answers $n$ questions by holding $n$ entries. The scheme developed
across Papers~1--2 answers them by holding a fixed table of twenty rows and
recomputing. In the phrasing of \citet{sachikonye2025c}: the dictionary is
empty; the laboratory is not.

\subsection{A live inconsistency in the corpus}

There is a second reason this paper is necessary, and it is not
architectural but corrective.

Paper~2 of this monograph---the CYP3A4 address manifold---invokes the Empty
Dictionary Principle by name, citing \citet{sachikonye2025c}, and describes
its own construction as ``the operational form'' of that principle. It then
writes, three paragraphs later:

\begin{quote}
\small
``The Empty Dictionary stores nothing; the entire content of
$\Sexp_{\mathrm{P450}}^{\mathrm{family}}$ is the syntactic concatenation of
addresses. Storage cost is $O(N \cdot k)$ trits per sequence. For $N = 500$
residues and $k = 9$ trits per residue, this is $4500$ trits ${\approx}7$
kbits per sequence; for $4 \times 10^{5}$ sequences, ${\sim}350$~MB.''
\end{quote}

``Stores nothing'' is $O(1)$. ``${\sim}350$~MB for $4 \times 10^{5}$
sequences'' is $O(n)$. These are not the same claim, and the passage asserts
both in five lines. Section~\ref{sec:reconcile} resolves this. The
resolution is not that Paper~2 miscalculated---we reproduce both its figures
to within its own stated precision---but that it attributes to the
\emph{scheme} a cost belonging to an optional \emph{cache}.

\subsection{Organisation}

Part~I fixes the resident state and proves it constant. Part~II shows that
queries are answered without entries, and---this is the part that requires
care---that the answers carry physical content rather than arithmetic
ceremony. Part~III reconciles Paper~2. Part~IV reports the validation, which
is exhaustive rather than sampled, and each of whose experiments carries a
control designed to fail if the corresponding claim were empty.

%% ============================================================
\part{The Resident State}
%% ============================================================
\section{What must be kept}
\label{sec:resident}

\begin{definition}[Residue coordinate table]
\label{def:table}
Let $\Res$ be the twenty canonical amino acids. The coordinate table is the
map
\[
  \Scoord : \Res \longrightarrow [0,1]^{3},
  \qquad
  \Scoord(a) = (\Sk(a), \St(a), \Se(a)),
\]
assigning to each residue its kinetic, thermal, and electronic S-entropy
coordinates as fixed in Paper~1.
\end{definition}

\begin{definition}[Encoding rule]
\label{def:encoding}
For $x \in [0,1]$ let $\trit_{j}(x)$ denote the $j$-th digit of the base-$3$
expansion of $x$. The address of a residue at resolution $k = 3m$ is the
interleaving
\[
  \addr_{k}(a) \;=\;
  \bigl(\trit_{1}\Sk,\, \trit_{1}\St,\, \trit_{1}\Se,\;
        \trit_{2}\Sk,\, \dots,\, \trit_{m}\Se\bigr),
\]
evaluated at $\Scoord(a)$. The encoding rule is the triple (base $b=3$,
interleave order $(\Sk,\St,\Se)$, domain $[0,1]$).
\end{definition}

\begin{definition}[Scheme and cache]
\label{def:scheme-cache}
The \emph{scheme} $\Sch$ is the pair (coordinate table, encoding rule). A
\emph{cache} $\Cache$ is any set of materialised addresses of specific
objects. We say a query procedure is \emph{cache-free} if its output is a
function of $\Sch$ and the query alone.
\end{definition}

The distinction in Definition~\ref{def:scheme-cache} is the one Paper~2
elides, and everything in Part~III turns on it.

\begin{proposition}[Resident state]
\label{prop:resident}
$|\Sch|$ is independent of the number of addressable objects. It is
$20$ coordinate rows plus a fixed encoding rule.
\end{proposition}

\begin{proof}
By Definition~\ref{def:table}, $\Scoord$ has domain $\Res$ with
$|\Res| = 20$, a constant fixed by the genetic code and not by any corpus.
By Definition~\ref{def:encoding}, $\addr_{k}$ is determined by $\Scoord$
together with the encoding rule, which is three constants. No object
outside $\Res$ appears in either definition. Hence $\Sch$ is fixed before
any corpus is named, and $|\Sch| = O(1)$ in the number of addressable
objects.
\end{proof}

\begin{remark}[The constant is small, but that is not the point]
\label{rem:small}
Serialised, $\Sch$ measures \SI{561}{\byte}
(Table~\ref{tab:scaling}). The number is unremarkable; what matters is its
\emph{invariance}. A scheme occupying a gigabyte but not growing with the
corpus would make the same claim. Experiment~01 therefore tests constancy
across six corpora spanning nine orders of magnitude, not smallness.
\end{remark}

\begin{corollary}[Addressability without enumeration]
\label{cor:addressable}
Every sequence over $\Res$ is addressable under $\Sch$, including sequences
that have never been observed. The set of addressable objects is
$\Res^{*}$, which is countably infinite, while $|\Sch|$ is finite.
\end{corollary}

\begin{proof}
Immediate from Definition~\ref{def:encoding}: $\addr_{k}$ is total on
$\Res$, and its extension to sequences is positionwise. Nothing in the
construction consults a list of known sequences.
\end{proof}

Corollary~\ref{cor:addressable} is where the departure from a database
becomes sharp. A database's addressable set is its key set, and grows only
by insertion. Here the addressable set is fixed at $\Res^{*}$ from the
outset and cannot be grown, because it is already everything.

\paragraph{What Part~I has, and what it lacks.}
Part~I establishes that the resident state is constant, and that the
addressable set is unbounded. It establishes nothing whatever about whether
the addresses are \emph{useful}. A scheme mapping every residue to the same
constant would satisfy Proposition~\ref{prop:resident} and
Corollary~\ref{cor:addressable} exactly, while answering no question at all.
Part~II supplies what is missing.

%% ============================================================
\part{Query Without Entries}
%% ============================================================
\section{Answering unstored questions}
\label{sec:query}

\begin{definition}[The three queries]
\label{def:queries}
For sequences $u, v \in \Res^{*}$:
\begin{itemize}
  \item \textsc{Identify}$(u)$: return $\addr_{k}(u)$.
  \item \textsc{Similar}$(u,v)$: return
    $\sigma(u,v) = \frac{1}{nk}\sum_{i=1}^{n} \ell\bigl(\addr_{k}(u_{i}),
    \addr_{k}(v_{i})\bigr)$, where $\ell$ is longest common prefix length
    and $n = \min(|u|,|v|)$.
  \item \textsc{Predict}$(u; \mathcal{C})$: return
    $\arg\max_{c \in \mathcal{C}} \sigma(u,c)$ for a reference set
    $\mathcal{C}$.
\end{itemize}
\end{definition}

\begin{theorem}[Empty dictionary, P450 instance]
\label{thm:empty}
Each query of Definition~\ref{def:queries} is cache-free in the sense of
Definition~\ref{def:scheme-cache}, and its evaluation leaves $\Sch$
unmodified.
\end{theorem}

\begin{proof}
\textsc{Identify} is the positionwise application of $\addr_{k}$, which by
Definition~\ref{def:encoding} reads only $\Scoord$ and the encoding rule.
\textsc{Similar} is a function of two \textsc{Identify} outputs and $\ell$,
a pure function on trit strings. \textsc{Predict} is an $\arg\max$ over
finitely many \textsc{Similar} evaluations. Hence all three are functions of
$\Sch$ and their arguments alone. None contains a write to $\Scoord$ or to
the encoding rule, so $\Sch$ is unmodified.
\end{proof}

\begin{corollary}[No entry required]
\label{cor:no-entry}
The queries are answerable for $u$ never previously observed.
\end{corollary}

\begin{proof}
By Theorem~\ref{thm:empty} the output depends on $u$ only through
$\addr_{k}(u)$, which is total on $\Res^{*}$ by
Corollary~\ref{cor:addressable}.
\end{proof}

\subsection{Content, not ceremony}
\label{sec:content}

Theorem~\ref{thm:empty} is a statement about \emph{plumbing}. It says the
queries do not read a cache. It does not say the answers mean anything, and
here we must be careful, because a scheme can satisfy
Theorem~\ref{thm:empty} perfectly while computing nothing of interest---the
constant map again.

The substantive claim is that address proximity tracks physicochemical
proximity. Testing it requires a task that a \emph{wrong} table would fail.
Our first attempt used one that would not, and the failure is instructive
enough to record in the manuscript rather than the appendix.

\begin{remark}[A control that refused to fail]
\label{rem:control-refused}
Experiment~02 originally posed \textsc{Predict} over closely related P450
substrate-recognition-site fragments~\citep{Gotoh1992}: given a CYP3A5
fragment, find its nearest CYP3A4 relative. The real table scored $1.00$.
So did a \emph{shuffled} table---same twenty rows, permuted assignment,
wrong physics---also $1.00$, with empirical $p = 1.000$.

The claim was not thereby refuted; the task was defective. The fragments sit
at ${\approx}94\,\%$ identity, so the nearest neighbour is determined by
string identity alone. \emph{Any} injective residue$\to$address map ranks the
near-identical pair first, whether its coordinates encode hydrophobicity or
a telephone directory. A task decided by identity cannot discriminate
physics from bookkeeping.

We report this because a control that cannot fail is not a control, and a
reader has no way to tell the two apart from a table of successes. The
non-discriminating task is retained in the output under
\texttt{task\_identity\_NON\_DISCRIMINATING}, carrying its own $p = 1.000$.
\end{remark}

The replacement task is graded rather than categorical. Let $\pi(a,b)$ be a
reference physicochemical distance built from Kyte--Doolittle
hydropathy~\citep{KyteDoolittle1982}, side-chain
volume~\citep{Zamyatnin1972}, and formal charge---three scales external to
the scheme and used nowhere in it---and let
\[
  \delta(a,b) \;=\; 1 - \tfrac{1}{k}\,\ell\bigl(\addr_{k}(a), \addr_{k}(b)\bigr)
\]
be address separation.

\begin{proposition}[Graded correlation]
\label{prop:graded}
Over all $\binom{20}{2} = 190$ residue pairs, $\delta$ and $\pi$ are rank
correlated~\citep{Spearman1904} at $\rho = 0.382$. Under $200$ random
permutations of the coordinate table, mean $\rho = -0.004$ with maximum
$0.252$; no permutation attains the observed value, giving empirical
$p < 0.005$.
\end{proposition}

\begin{proof}[Measurement]
Experiment~02, exhaustive over pairs rather than sampled. See
Table~\ref{tab:query}.
\end{proof}

\begin{remark}[On the size of $\rho$]
\label{rem:rho-honest}
$\rho = 0.382$ is a moderate correlation, and we decline to describe it as
more. Two things bound it from above. The interleaved address is a
\emph{lexicographic} object: a residue pair differing in the first trit of
$\Sk$ shares no prefix at all, however close in volume, so $\delta$ is a
coarsened and discontinuous function of coordinate distance by construction.
And $\pi$ weights three scales equally, which no physical argument compels.
The claim supported is that the coordinate assignment carries
physicochemical information no permutation of it reproduces---the separation
from the null is unambiguous---not that $\delta$ is a good metric for
physicochemical distance. It is not.
\end{remark}

\paragraph{What Part~II has, and what it lacks.}
Part~II establishes that queries are answered without entries
(Theorem~\ref{thm:empty}), and that the answers depend on the specific
coordinate assignment rather than on any injective relabelling
(Proposition~\ref{prop:graded}). It does not establish that the scheme
answers \emph{as well as} a database would. On tasks where a database holds
the exact answer, retrieval wins whenever the object was stored;
Corollary~\ref{cor:no-entry} concerns the complementary case, and the
comparison in Table~\ref{tab:query} is confined to it.

%% ============================================================
\part{Reconciliation with Paper~2}
%% ============================================================
\section{Which quantity is \texorpdfstring{$350$}{350}~MB?}
\label{sec:reconcile}

We first confirm that no arithmetic is in dispute.

\begin{proposition}[Paper~2's figures reproduce]
\label{prop:p2-arith}
With $N = 500$, $k = 9$, and $4 \times 10^{5}$ sequences, packing trits at
the information-theoretic limit $\log_{2} 3$ bits each gives $7.13$~kbits
per sequence and $356.6$~MB for the corpus, against Paper~2's stated
${\approx}7$~kbits and ${\sim}350$~MB.
\end{proposition}

\begin{proof}[Measurement]
Experiment~03. Both agree to within Paper~2's stated precision, so the
disagreement below is not arithmetical.
\end{proof}

\begin{theorem}[Attribution]
\label{thm:attribution}
Paper~2's $O(N \cdot k)$ figure measures $|\Cache|$ for a corpus cache, not
$|\Sch|$. Under Definition~\ref{def:scheme-cache} these are different
objects, and only the latter is the resident state.
\end{theorem}

\begin{proof}
By Proposition~\ref{prop:resident}, $|\Sch| = O(1)$ and is fixed before any
corpus is named; the quantity $4500$ trits is indexed by a \emph{sequence},
hence cannot be a measurement of $\Sch$. It is by construction the size of
the materialised address of one object, i.e.\ a cache entry. Summing over
$4 \times 10^{5}$ objects yields $|\Cache|$, which is $O(n)$.
\end{proof}

\begin{proposition}[The cache is not load-bearing]
\label{prop:cache-optional}
Let $A_{\Cache}$ and $A_{\emptyset}$ be the answers to the full query set
with the cache materialised and deleted respectively. Then
$A_{\Cache} = A_{\emptyset}$ byte for byte, and $|\Sch|$ is unchanged by
either.
\end{proposition}

\begin{proof}[Measurement]
Experiment~03, which runs both and compares serialisations. The equality is
forced by Theorem~\ref{thm:empty}---no query reads $\Cache$---but is
measured rather than assumed, because the theorem constrains the
\emph{definition} of the queries and not their implementation.
\end{proof}

\begin{remark}[What Paper~2 should say]
\label{rem:p2-fix}
Paper~2's sentence ``The Empty Dictionary stores nothing $\dots$ storage cost
is $O(N \cdot k)$ trits per sequence'' conflates the two. The corrected
statement separates them:

\begin{quote}
\small
The scheme stores nothing about any sequence: its resident state is the
twenty-row coordinate table and the encoding rule, \SI{561}{\byte},
independent of corpus size. \emph{Should one choose} to materialise
addresses for a fixed corpus---an optional cache, not a precondition for
querying---that cache costs $O(N \cdot k)$ trits per sequence,
${\sim}350$~MB for $4 \times 10^{5}$ sequences, five orders of magnitude
below the corresponding structural
representation~\citep{Berman2000,Varadi2022}.
\end{quote}

The compression claim against PDB and AlphaFold-DB survives intact; it was
always a claim about the cache, and it is a good one. What does not survive
is calling that cache an empty dictionary. A $350$~MB artefact whose size is
proportional to the corpus is a \emph{compressed} dictionary. The
distinction is not pedantic: a compressed dictionary cannot answer for an
object outside it, and Corollary~\ref{cor:no-entry} can.
\end{remark}

\begin{remark}[The honest form of the claim]
\label{rem:empty-honest}
Stated without inflation: storage is $O(1)$ where a database is $O(n)$, and
this is bought by recomputation, not by cleverness about compression. Every
query pays $O(Nk)$ arithmetic it would not pay against an index. For a
corpus that is fixed, small, and queried often, the database is simply
better, and nothing here argues otherwise. The regime where the empty scheme
wins is the one where the query set is not known in advance---which, for the
$>300$ PharmVar alleles~\citep{Gaedigk2018} and the open-ended space of
engineered P450 variants, is the usual regime.

And the coordinates themselves came from measurement. Twenty rows had to be
determined experimentally before any of this could run. The dictionary is
empty; the laboratory is not.
\end{remark}

\paragraph{What Part~III has, and what it lacks.}
Part~III attributes each of Paper~2's two sentences to a distinct quantity
and shows by measurement that the larger one is optional. It does not
propose an erratum procedure for the monograph, and it does not re-derive
Paper~2's compression comparison against PDB, which we accept as stated.

%% ============================================================
\part{Validation}
%% ============================================================
\section{Design}
\label{sec:validation}

Three experiments. Each is exhaustive over its domain rather than sampled,
and each carries a control constructed so that it would \emph{fail} if the
claim under test were empty. Controls that cannot fail are worth nothing,
which is the lesson of Remark~\ref{rem:control-refused}.

\begin{table}[H]
\centering
\small
\caption{Experiment 01 --- resident state against corpus size. The resident
column is the claim; the cache column is the control, which must grow.}
\label{tab:scaling}
\begin{tabular}{lrrr}
\toprule
Corpus & $n$ & Cache (MB) & $\Sch$ (B) \\
\midrule
Single CYP3A4      & $1$                & $0.001$      & $561$ \\
Human isoforms     & $57$               & $0.051$      & $561$ \\
PharmVar alleles   & $300$              & $0.267$      & $561$ \\
UniProt P450 set   & $4\times10^{5}$    & $356.6$      & $561$ \\
Hypothetical       & $10^{7}$           & $8915.4$     & $561$ \\
Hypothetical       & $10^{9}$           & $891541.4$   & $561$ \\
\bottomrule
\end{tabular}
\end{table}

The resident column is constant across nine orders of magnitude; the control
column spans six. Had both been constant, the comparison would have been
empty.

\begin{table}[H]
\centering
\small
\caption{Experiment 02 --- query without entries, with two controls.}
\label{tab:query}
\begin{tabular}{lr}
\toprule
Measurement & Value \\
\midrule
Queries answered, empty scheme        & $6/6$ \\
Queries answered, control A (index)   & $0/6$ \\
$\Sch$ unchanged by query batch       & yes \\
\midrule
Graded task, real table $\rho$        & $0.382$ \\
Control B, mean $\rho$ ($200$ shuffles) & $-0.004$ \\
Control B, max $\rho$                 & $0.252$ \\
Empirical $p$                         & $<0.005$ \\
Residue pairs                         & $190$ \\
\midrule
Identity task, real table             & $1.00$ \\
Identity task, shuffled mean          & $1.00$ \\
Identity task, empirical $p$          & $1.000$ \\
\bottomrule
\end{tabular}
\end{table}

The last three rows are the non-discriminating task of
Remark~\ref{rem:control-refused}, reported rather than discarded. Their
$p = 1.000$ is the signature of a task that tests nothing, and it is worth
having on the page next to the $p < 0.005$ of one that does.

\begin{table}[H]
\centering
\small
\caption{Experiment 03 --- reconciliation with Paper~2.}
\label{tab:reconcile}
\begin{tabular}{lr}
\toprule
Measurement & Value \\
\midrule
Paper~2 per sequence, recomputed  & $7.13$ kbits \\
Paper~2 per sequence, as stated   & ${\approx}7$ kbits \\
Paper~2 corpus, recomputed        & $356.6$ MB \\
Paper~2 corpus, as stated         & ${\sim}350$ MB \\
Resident state $|\Sch|$           & \SI{561}{\byte} \\
Answers with vs.\ without cache   & identical \\
$|\Sch|$ affected by cache        & no \\
\bottomrule
\end{tabular}
\end{table}

\begin{remark}[Scope of the validation]
\label{rem:validation-scope}
Three limits, stated plainly. The query corpus is ten P450 fragments, not a
proteome: it is large enough to separate the empty scheme from an index
($6/6$ against $0/6$) and no larger, and the answerability claim is
structural rather than statistical. The graded task uses three
physicochemical scales chosen before the correlation was computed but not
pre-registered. And every experiment tests the scheme against \emph{itself}
under permutation, not against an external method; nothing here shows the
addressing is better than BLAST at anything, and no such claim is made.
\end{remark}

\paragraph{What Part~IV has, and what it lacks.}
Part~IV supplies measured values for every number asserted in Parts~I--III,
with controls that failed informatively in one case and were replaced. It
supplies no comparison against established sequence-analysis tooling, and
its corpus is small.

%% ============================================================
\section{Conclusion}
%% ============================================================

The seventeenth paper of this monograph asks the question the other sixteen
presuppose away: whether P450 knowledge must be stored. It need not. The
resident state of the addressing scheme is twenty rows and an encoding rule,
\SI{561}{\byte}, constant from one sequence to $10^{9}$, and the three
queries of Definition~\ref{def:queries} are answered for objects that were
never entered---$6/6$ where an exact-match index answers $0/6$---with the
resident state provably and measurably unmodified.

The answers carry physical content, and we established this only after a
first control refused to fail. That episode is left in the manuscript
deliberately. The graded correlation $\rho = 0.382$, which no permutation out
of $200$ reproduces, is a modest number honestly obtained; the perfect $1.00$
it replaced was a strong number that meant nothing.

Finally, Paper~2's storage passage is corrected. Its arithmetic is right and
reproduces to its stated precision; its attribution is wrong. The
${\sim}350$~MB it reports is an optional cache, demonstrably not
load-bearing, and not the resident state. Paper~2 describes a compressed
dictionary and calls it an empty one. The compression claim against
structural databases stands; the identification with the Empty Dictionary
Principle does not, because a compressed dictionary cannot answer for what
it does not contain, and this scheme can.

\bibliographystyle{plainnat}
\bibliography{references}

\end{document}
